% RQ2: anatomy of the eight unsafe fresh-model admits.
%
% Built per section 5.4 of docs/rq1-rq3-metrics-guide.md from the workbooks
% below (ML implementation commit c871d3a); same decision set as the
% unsafe-admit cell of Table~\ref{tab:rq2_admission_audit}.  Rows are ordered
% by spike size within each optional-work group, so the labels are positional.
%
%   A = 48U_metrics_12MP_normal_0729_PacingOnly_1.xlsx
%   B = 48U_metrics_24MP_memory_0729_PacingOnly_1.xlsx
%   C = 48U_metrics_24MP_memory_0729_PacingOnly_2.xlsx
%
%   id  book  runId  shot     B     Cbar      C      UB   C/Cbar  overrun
%   M1    A     30     21   1516    861.7   2033   903.8   2.36      517
%   M2    A     11      9   1308   1123.7   2367  1229.1   2.11     1059
%   M3    A      9     26   1749   1189.3   2357  1222.5   1.98      608
%   M4    A     31     29    746    607.0    777   682.8   1.28       31
%   M5    C     18     20    844    684.7    855   765.1   1.25       11
%   S1    C      6     11    975    625.7    979   915.0   1.56        4
%   S2    B     29     14    653    544.0    713   593.0   1.31       60
%   S3    C     18     20    472    397.3    483   467.7   1.22       11
%
% Panel (b) normalized factors.  Over the remaining node suffix, camera-process
% CPU time is sum(cpuTimeMs), wall time is sum(wallTimeMs), and average busy
% cores is CPU time / wall time.  Each event is divided by the unweighted mean
% of the same quantity over its three recent-safe decisions.
% Values, CPU-time factor / average-busy-cores factor:
%   M1  1.443 / 0.602     M5  1.272 / 1.011
%   M2  1.346 / 0.634     S1  1.055 / 0.669
%   M3  1.501 / 0.750     S2  1.532 / 1.176
%   M4  1.100 / 0.859     S3  1.261 / 1.031
% The corresponding wall-time factor is CPU-time factor / busy-cores factor.
% Blocking GC time was 0 ms in every baseline and every event; overheat level
% and thermal status were identical to the recent-safe baseline in all eight.

% Layout: one panel per row, each spanning the column, so the eight rows get
% the full width instead of half of it.
\begingroup
\def\rqtwoReviewerAccounting{1}

\begin{figure}[!t]
  \centering
  % Styles and legend labels are defined before the picture so the panel
  % options below can read them; each panel then draws its own key inside its
  % axis.
    \tikzset{
      glabel/.style={font=\tiny,inner sep=1pt},
    }
    \pgfplotsset{
      spikepanel/.style={
        ymin=0.35,
        ymax=8.65,
        ytick={1,2,3,4,5,6,7,8},
        yticklabels={S3,S2,S1,M5,M4,M3,M2,M1},
        axis x line*=bottom,
        axis y line*=left,
        % Outward ticks, as in fig_rq3_pacing_trajectories: these rows are
        % dense, and inward ticks both intrude on the markers and duplicate
        % the x grid lines.
        tick align=outside,
        xmajorgrids,
        grid style={black!14,line width=0.3pt},
        axis line style={line width=0.35pt},
        tick style={line width=0.35pt},
        clip=false,
        legend cell align=left,
      },
    }
    \pgfplotsset{
      spikepanel/.append style={
        % Sized so one panel plus its labels fills \columnwidth; the
        % \resizebox below then runs at a scale of about 1.0.
        width=1.105\columnwidth,
        % 4.50cm over the 8.3-unit range gives a row pitch of about 15.5pt,
        % roughly three times the tallest marker, so the rows read apart.
        height=4.50cm,
        tick label style={font=\tiny},
        yticklabel style={font=\tiny},
        label style={font=\tiny},
        % Panel tags sit above their panel, as in
        % fig_rq3_pacing_trajectories: stacked vertically, a tag below panel
        % (a) reads as if it belonged to panel (b).
        title style={font=\tiny},
        xlabel style={yshift=2pt},
        % In-axis key: at full column width the Single-frame rows leave the
        % right of that band empty in both panels.  Each panel sets its own
        % anchor point, flush right and at y=1.925, the centre of the band
        % between the multi-frame separator (3.5) and ymin (0.35).  Anchoring
        % on east rather than south east centres the box there whatever its
        % height, so the two keys balance despite holding three entries and
        % two.  Only the bottom and left axis lines are drawn, so nothing
        % collides with the right edge.
        legend style={
          font=\tiny,
          draw=black!25,
          fill=white,
          line width=0.2pt,
          inner sep=1pt,
          row sep=-1pt,
          legend columns=1,
          anchor=east,
        },
      },
    }
    \ifdefined\rqtwoReviewerAccounting
      \def\rqtwoBaselineLegend{Recent-safe avg.}
      \def\rqtwoBoundLegend{Predicted bound}
      \def\rqtwoEventLegend{Unsafe admit}
    \else
      \def\rqtwoBaselineLegend{Baseline}
      \def\rqtwoBoundLegend{Bound}
      \def\rqtwoEventLegend{Event}
    \fi
    % Panel (b) shows the two measured factors in
    % wall-time factor = CPU-time factor / average-busy-cores factor.
    \def\rqtwoCpuLegend{CPU time}
    \def\rqtwoParallelLegend{Average busy cores}

  \newsavebox{\rqtwospikebox}
  \begin{lrbox}{\rqtwospikebox}
  \begin{tikzpicture}
    \begin{groupplot}[
      group style={
        group size=1 by 2,
        % Room for the lower edge of panel (a), tick labels plus a one-line
        % x label, and for panel (b)'s tag above it: two text lines have to
        % fit between the plot boxes without touching.
        vertical sep=1.60cm,
      },
      spikepanel,
    ]

    %% ---------------------------------------------------------------- (a)
      \nextgroupplot[
        % Headroom past 2.0 keeps the growth-factor labels inside the axis.
        xmin=0.46,
        xmax=2.05,
        xtick={0.5,1.0,1.5,2.0},
        xticklabels={0.5,1.0,1.5,2.0},
        title={(a) Remaining latency vs. budget},
        xlabel={Remaining-sequence latency $/$ admission budget},
        legend style={at={(axis cs:2.05,1.925)}},
      ]
      % Admission budget and the multi-frame / single-frame separator.  Both
      % rules span the full panel: the key sits inside the Single-frame band,
      % clear of the y=3.5 separator, and its white fill covers the grid.
      \draw[black!60,densely dashed,line width=0.5pt]
        (axis cs:1,0.35) -- (axis cs:1,8.65);
      \node[glabel,anchor=south,text=black!65] at (axis cs:1,8.65) {budget};
      \draw[black!70,line width=0.6pt] (axis cs:0.46,3.5) -- (axis cs:2.05,3.5);

      % Recent-safe baseline -> unsafe event.
      \addplot[black!45,->,>=stealth,line width=0.6pt,forget plot]
        coordinates {(0.568,8) (1.341,8)};
      \addplot[black!45,->,>=stealth,line width=0.6pt,forget plot]
        coordinates {(0.859,7) (1.810,7)};
      \addplot[black!45,->,>=stealth,line width=0.6pt,forget plot]
        coordinates {(0.680,6) (1.348,6)};
      \addplot[black!45,->,>=stealth,line width=0.6pt,forget plot]
        coordinates {(0.814,5) (1.042,5)};
      \addplot[black!45,->,>=stealth,line width=0.6pt,forget plot]
        coordinates {(0.811,4) (1.013,4)};
      \addplot[black!45,->,>=stealth,line width=0.6pt,forget plot]
        coordinates {(0.642,3) (1.004,3)};
      \addplot[black!45,->,>=stealth,line width=0.6pt,forget plot]
        coordinates {(0.833,2) (1.092,2)};
      \addplot[black!45,->,>=stealth,line width=0.6pt,forget plot]
        coordinates {(0.842,1) (1.023,1)};

      \addplot[only marks,mark=o,mark size=2.0pt,draw=black,fill=white,
               line width=0.5pt]
        coordinates {(0.568,8) (0.859,7) (0.680,6) (0.814,5)
                     (0.811,4) (0.642,3) (0.833,2) (0.842,1)};
      \addlegendentry{\rqtwoBaselineLegend}

      \addplot[only marks,mark=|,mark size=2.7pt,draw=black!60,
               line width=0.7pt]
        coordinates {(0.596,8) (0.940,7) (0.699,6) (0.915,5)
                     (0.907,4) (0.938,3) (0.908,2) (0.991,1)};
      \addlegendentry{\rqtwoBoundLegend}

      \addplot[only marks,mark=square*,mark size=1.9pt,draw=black,fill=black]
        coordinates {(1.341,8) (1.810,7) (1.348,6) (1.042,5)
                     (1.013,4) (1.004,3) (1.092,2) (1.023,1)};
      \addlegendentry{\rqtwoEventLegend}

      % Growth over the recent-safe baseline.
      \node[glabel,anchor=west,xshift=2.5pt] at (axis cs:1.341,8) {$\times$2.36};
      \node[glabel,anchor=west,xshift=2.5pt] at (axis cs:1.810,7) {$\times$2.11};
      \node[glabel,anchor=west,xshift=2.5pt] at (axis cs:1.348,6) {$\times$1.98};
      \node[glabel,anchor=west,xshift=2.5pt] at (axis cs:1.042,5) {$\times$1.28};
      \node[glabel,anchor=west,xshift=2.5pt] at (axis cs:1.013,4) {$\times$1.25};
      \node[glabel,anchor=west,xshift=2.5pt] at (axis cs:1.004,3) {$\times$1.56};
      \node[glabel,anchor=west,xshift=2.5pt] at (axis cs:1.092,2) {$\times$1.31};
      \node[glabel,anchor=west,xshift=2.5pt] at (axis cs:1.023,1) {$\times$1.22};

    %% ---------------------------------------------------------------- (b)
      \nextgroupplot[
        xmin=0.50,
        xmax=1.62,
        xtick={0.6,0.8,1.0,1.2,1.4,1.6},
        xticklabels={0.6,0.8,1.0,1.2,1.4,1.6},
        title={(b) CPU time and parallelism changes},
        xlabel={Relative to recent-safe average ($\times$)},
        legend style={at={(axis cs:1.61,2.88)}},
      ]
      \addplot[draw=none,forget plot] coordinates {(1,1)};

      \addlegendimage{only marks,mark=square*,mark size=1.9pt,
                      draw=black,fill=black}
      \addlegendentry{\rqtwoCpuLegend}
      \addlegendimage{only marks,mark=o,mark size=2.0pt,
                      draw=black,fill=white,line width=0.5pt}
      \addlegendentry{\rqtwoParallelLegend}

      % A factor of 1.0 is the recent-safe average.  Offsetting the two series
      % within each row makes their directions readable without encoding raw
      % metric changes as same-direction latency contributions.
      \draw[black!60,densely dashed,line width=0.5pt]
        (axis cs:1,0.35) -- (axis cs:1,8.65);
      \node[glabel,anchor=south,text=black!65]
        at (axis cs:1,8.65) {recent-safe avg.};
      \draw[black!70,line width=0.6pt]
        (axis cs:0.50,3.5) -- (axis cs:1.62,3.5);

      % Change segments from the normalized recent-safe value.
      % CPU time (upper position in each row).
      \draw[black!45,line width=0.6pt] (axis cs:1,8.17) -- (axis cs:1.443,8.17);
      \draw[black!45,line width=0.6pt] (axis cs:1,7.17) -- (axis cs:1.346,7.17);
      \draw[black!45,line width=0.6pt] (axis cs:1,6.17) -- (axis cs:1.501,6.17);
      \draw[black!45,line width=0.6pt] (axis cs:1,5.17) -- (axis cs:1.100,5.17);
      \draw[black!45,line width=0.6pt] (axis cs:1,4.17) -- (axis cs:1.272,4.17);
      \draw[black!45,line width=0.6pt] (axis cs:1,3.17) -- (axis cs:1.055,3.17);
      \draw[black!45,line width=0.6pt] (axis cs:1,2.17) -- (axis cs:1.532,2.17);
      \draw[black!45,line width=0.6pt] (axis cs:1,1.17) -- (axis cs:1.261,1.17);

      % Average busy cores (lower position in each row).
      \draw[black!45,line width=0.6pt] (axis cs:1,7.83) -- (axis cs:0.602,7.83);
      \draw[black!45,line width=0.6pt] (axis cs:1,6.83) -- (axis cs:0.634,6.83);
      \draw[black!45,line width=0.6pt] (axis cs:1,5.83) -- (axis cs:0.750,5.83);
      \draw[black!45,line width=0.6pt] (axis cs:1,4.83) -- (axis cs:0.859,4.83);
      \draw[black!45,line width=0.6pt] (axis cs:1,3.83) -- (axis cs:1.011,3.83);
      \draw[black!45,line width=0.6pt] (axis cs:1,2.83) -- (axis cs:0.669,2.83);
      \draw[black!45,line width=0.6pt] (axis cs:1,1.83) -- (axis cs:1.176,1.83);
      \draw[black!45,line width=0.6pt] (axis cs:1,0.83) -- (axis cs:1.031,0.83);

      \addplot[only marks,mark=square*,mark size=1.9pt,
               draw=black,fill=black]
        coordinates {(1.443,8.17) (1.346,7.17) (1.501,6.17) (1.100,5.17)
                     (1.272,4.17) (1.055,3.17) (1.532,2.17) (1.261,1.17)};
      \addplot[only marks,mark=o,mark size=2.0pt,
               draw=black,fill=white,line width=0.5pt]
        coordinates {(0.602,7.83) (0.634,6.83) (0.750,5.83) (0.859,4.83)
                     (1.011,3.83) (0.669,2.83) (1.176,1.83) (1.031,0.83)};

      % Direct factor labels; labels sit outside their change segment.
      \node[glabel,anchor=west,xshift=2.5pt] at (axis cs:1.443,8.17) {$1.44\times$};
      \node[glabel,anchor=east,xshift=-2.5pt] at (axis cs:0.602,7.83) {$0.60\times$};
      \node[glabel,anchor=west,xshift=2.5pt] at (axis cs:1.346,7.17) {$1.35\times$};
      \node[glabel,anchor=east,xshift=-2.5pt] at (axis cs:0.634,6.83) {$0.63\times$};
      \node[glabel,anchor=west,xshift=2.5pt] at (axis cs:1.501,6.17) {$1.50\times$};
      \node[glabel,anchor=east,xshift=-2.5pt] at (axis cs:0.750,5.83) {$0.75\times$};
      \node[glabel,anchor=west,xshift=2.5pt] at (axis cs:1.100,5.17) {$1.10\times$};
      \node[glabel,anchor=east,xshift=-2.5pt] at (axis cs:0.859,4.83) {$0.86\times$};
      \node[glabel,anchor=west,xshift=2.5pt] at (axis cs:1.272,4.17) {$1.27\times$};
      \node[glabel,anchor=west,xshift=2.5pt] at (axis cs:1.011,3.83) {$1.01\times$};
      \node[glabel,anchor=west,xshift=2.5pt] at (axis cs:1.055,3.17) {$1.05\times$};
      \node[glabel,anchor=east,xshift=-2.5pt] at (axis cs:0.669,2.83) {$0.67\times$};
      \node[glabel,anchor=west,xshift=2.5pt] at (axis cs:1.532,2.17) {$1.53\times$};
      \node[glabel,anchor=west,xshift=2.5pt] at (axis cs:1.176,1.83) {$1.18\times$};
      \node[glabel,anchor=west,xshift=2.5pt] at (axis cs:1.261,1.17) {$1.26\times$};
      \node[glabel,anchor=west,xshift=2.5pt] at (axis cs:1.031,0.83) {$1.03\times$};

    \end{groupplot}
  \end{tikzpicture}
  \end{lrbox}
  \resizebox{\columnwidth}{!}{\usebox{\rqtwospikebox}}

  \caption{Anatomy of the eight unsafe admits the always-admit model would
  have produced.  (a) Remaining-sequence latency normalized by the admission
  budget, with the growth over the recent-safe average annotated per row.
  (b) The camera process's CPU time and average busy cores over the remaining
  sequence, each normalized by its recent-safe average; 1.0 denotes no change.
  CPU time rises for every unsafe admit, whereas average busy cores falls for
  M1--M4 and S1 and rises for M5, S2, and S3.}
  \label{fig:rq2_unsafe_spike_anatomy}

\end{figure}
\endgroup
